% !TeX root = ../../../../../main.tex
% cs.tex

\subsection{CS（消費供給）曲線の導出}
\label{sec:chapter5_asad_cs}
CS曲線は、供給側の式をすべて集約したAS曲線を出発点として導出される。

AS曲線は代表的家計の価格最適化に関する全方程式（最適価格決定式、補助変数の動学、PPIの動学）によって構成され、価格硬直性が高い短期的状況下では生産者物価指数（PPI）\(p_t^H\)を一定水準に固定する供給制約として機能する。
このAS曲線によって確定した供給側の制約\(p_t^H\)を、消費平面\((c_t^{H\to W},p_t^{H\to W})\)へと「翻訳」するために、CPIの定義式を用いる。
\begin{equation}
\label{eq:cs_cpi_definition}
    p_t^{H\to W}=(p_t^H)^{\alpha^H}(e_t^{/*}p_t^{F*})^{1-\alpha^H}
\end{equation}
ここに、名目為替レートの決定式\(e_t^{/*}=\operatorname{E}_t[\mathcal{B}_t\Lambda_{\infty}](p_t^{H\to W}c_t^{H\to W})/(p_t^{F\to W*}c_t^{F\to W})\)を代入する。
名目為替レートの決定式は国際金融市場における資金の需給均衡（リスク共有条件）を表すものであり、特定の財の需要や供給に固有の方程式ではない。一般均衡モデルにおいてはシステム全体を繋ぐ変換ルールとして機能するため、どの変数を消去しどの変数を軸として残すかという分析目的に応じて、需要側の定式化にも本節のような供給側の定式化にも数学的に正当に組み込むことができる。この式を代入することで、国内の生産価値を消費価値に変換するための外部環境（交易条件の制約）が供給側に組み込まれる。
\begin{equation}
\label{eq:cs_cpi_substitution}
    p_t^{H\to W}=(p_t^H)^{\alpha^H}\left(\operatorname{E}_t[\mathcal{B}_t\Lambda_{\infty}]\frac{p_t^{H\to W}c_t^{H\to W}}{p_t^{F\to W*}c_t^{F\to W}}p_t^{F*}\right)^{1-\alpha^H}
\end{equation}
ここで、外国の変数と割引因子からなる自国政策から独立した項（遮断効果）を\(\Phi_t\equiv\operatorname{E}_t[\mathcal{B}_t\Lambda_{\infty}]p_t^{F*}/(p_t^{F\to W*}c_t^{F\to W})\)と定義して式を整理する。
\begin{equation}
\label{eq:cs_cpi_simplify}
    p_t^{H\to W}=(p_t^H)^{\alpha^H}\left(\Phi_tp_t^{H\to W}c_t^{H\to W}\right)^{1-\alpha^H}
\end{equation}
両辺を\((p_t^{H\to W})^{1-\alpha^H}\)で割り、\(\alpha^H\)乗根をとる（\(1/\alpha^H\)乗する）ことで、以下のCS曲線が得られる。
\begin{equation}
\label{eq:cs_curve}
    p_t^{H\to W}=p_t^H\Phi_t^{\frac{1-\alpha^H}{\alpha^H}}(c_t^{H\to W})^{\frac{1-\alpha^H}{\alpha^H}}
\end{equation}
式\eqref{eq:cs_curve}の指数\((1-\alpha^H)/\alpha^H\)が正の値をとるため、CS曲線は右上がりの曲線となる。この式が示す通り、CS曲線はAS曲線によって水平に固定されたPPI\(p_t^H\)を出発点（アンカー）として機能している。これは、実質消費\(c_t^{H\to W}\)が増加すると、名目為替レートの減価を通じて輸入品価格が上昇し、結果としてCPI\(p_t^{H\to W}\)が押し上げられるという供給側のトレードオフを完璧に数式化している。